%! Piecewise \& Step Functions — Unit 2, Lesson 2.4 — Homework (Answer Key)
\documentclass[10pt]{article}
\usepackage{algebra2-article}
\usepackage{algebra2-key}   % pulls in -boxes; do NOT also load -boxes
\newcommand{\TallMath}[1]{$\displaystyle #1\rule[-1.4em]{0pt}{3.2em}$}
% Coordinate grid: {xmin}{xmax}{ymin}{ymax}
\newcommand{\lgrid}[4]{%
  \draw[step=1, linegray!40, very thin] (#1,#3) grid (#2,#4);
  \draw[->, semithick, charcoal] (#1,0)--(#2,0) node[below right, font=\scriptsize]{$x$};
  \draw[->, semithick, charcoal] (0,#3)--(0,#4) node[left, font=\scriptsize]{$y$};}
% Endpoint dots
\newcommand{\closeddot}[2]{\fill[forest] (#1,#2) circle (3.5pt);}
\newcommand{\opendot}[2]{\draw[very thick, forest, fill=white] (#1,#2) circle (3.5pt);}

\begin{document}

\pageheader{Unit 2, Lesson 2.4}{Homework --- Answer Key}

\vspace{0.10in}

\begin{notesbox}{Practice}
Show your work. To evaluate a piecewise function, first find the piece whose condition the input
satisfies (watch the boundary), then apply that rule. On a graph, a \emph{closed} dot includes its
endpoint; an \emph{open} dot excludes it.

\begin{enumerate}[leftmargin=1.9em, itemsep=11pt]
  \item \textbf{Evaluate} $f(x)=\begin{cases} x+4, & x\le -1 \\ 2x, & x>-1 \end{cases}$
        \\[3pt]
        $f(-3)={\color{keyred}\mathbf{1}}$ \quad $f(-1)={\color{keyred}\mathbf{3}}$ \quad
        $f(0)={\color{keyred}\mathbf{0}}$ \quad $f(5)={\color{keyred}\mathbf{10}}$

  \item \textbf{Read the graph} of $f(x)=\begin{cases} x+2, & x<1 \\ 5-x, & x\ge 1 \end{cases}$.
        \begin{center}
        \begin{tikzpicture}[scale=0.5]
          \lgrid{-3}{5}{-1}{5}
          \draw[navy, dashed, semithick] (1,-1)--(1,5);
          \draw[very thick, forest] (-3,-1)--(1,3);
          \opendot{1}{3}
          \draw[very thick, forest] (1,4)--(4,1);
          \closeddot{1}{4}
        \end{tikzpicture}
        \end{center}
        $f(-1)={\color{keyred}\mathbf{1}}$; \; $f(1)={\color{keyred}\mathbf{4}}$; \; domain
        \ans{all real numbers}; \; increasing on \ans{$(-\infty,1)$}; \; decreasing on
        \ans{$[1,\infty)$}. \; Is $f$ continuous at $x=1$? \ans{no}.

  \item \textbf{Write it as pieces.} Rewrite $|x+1|$ as a two-piece function.
        \[ |x+1|=\begin{cases} {\color{keyred}\mathbf{x+1}}, & x\ge -1 \\[3pt] {\color{keyred}\mathbf{-x-1}}, & x<-1 \end{cases} \]

  \item \textbf{Greatest-integer function.} Evaluate.
        \\[2pt]
        $\lfloor 4.2\rfloor={\color{keyred}\mathbf{4}}$ \quad $\lfloor -3.1\rfloor={\color{keyred}\mathbf{-4}}$ \quad
        $\lfloor 6\rfloor={\color{keyred}\mathbf{6}}$ \quad $\lfloor -0.5\rfloor={\color{keyred}\mathbf{-1}}$. \\[2pt]
        For every $x$ in $[3,4)$, $\lfloor x\rfloor={\color{keyred}\mathbf{3}}$.

  \item \textbf{Model it.} A phone plan costs \$30 for up to $5$ GB of data, then \$10 for each GB over
        $5$. Let $C(g)$ be the monthly cost for $g$ GB.
        \begin{enumerate}[label=(\alph*), leftmargin=1.8em, itemsep=4pt]
          \item Complete the rule:
                $C(g)=\begin{cases} {\color{keyred}\mathbf{30}}, & 0\le g\le 5 \\[3pt] 30+{\color{keyred}\mathbf{10(g-5)}}, & g>5 \end{cases}$
          \item Find $C(3)$ and $C(8)$, and say what $C(8)$ means.
                \begin{work}
                  C(3) &= 30 \\
                  C(8) &= 30+10(8-5) \\
                       &= 30+10(3) \\
                       &= 60
                \end{work}
                \ansline{\$60 is the cost for 8 GB --- \$30 base plus 3 GB over 5 at \$10 each.}
        \end{enumerate}

  \item \textbf{Explain.} When you evaluate a piecewise function at $x=a$, how do you decide which piece
        to use? And how can you tell whether the function is \emph{continuous} at a boundary?
        \ansline{Use the piece whose condition $x=a$ satisfies (the $\le/\ge$ tells you which piece owns a boundary point).}
        \ansline{It is continuous at a boundary when both pieces give the \emph{same value} there; if they differ, the graph jumps.}
\end{enumerate}
\end{notesbox}

\begin{extensionbox}
\textbf{Build and reason.}
\begin{enumerate}[label=(\alph*), leftmargin=1.9em, itemsep=6pt]
  \item Write the piecewise rule for the function that is the constant $0$ for $x<0$, the line $y=2x$ for
        $0\le x\le 2$, and the constant $4$ for $x>2$. Is it continuous at $x=0$ and at $x=2$? Explain.
  \item A simple tax charges $10\%$ on income up to \$10{,}000 and $20\%$ on the amount above \$10{,}000.
        Writing $T(x)$ for the tax on income $x$, complete $T(x)=\{\,0.10x \text{ if } 0\le x\le 10000;\
        1000+{\color{keyred}\mathbf{0.20(x-10000)}} \text{ if } x>10000\,\}$, and explain why $T$ is
        continuous at \$10{,}000.
\end{enumerate}
\ansline{(a) $f(x)=\{\,0 \text{ if } x<0;\ 2x \text{ if } 0\le x\le 2;\ 4 \text{ if } x>2\,\}$; continuous at both boundaries.}
\ansline{At $x=0$ both pieces give $0$; at $x=2$ both give $4$ --- the pieces meet, so no jump.}
\ansline{(b) $T$ is continuous at \$10{,}000 because $0.10(10000)=1000$ matches $1000+0.20(0)=1000$.}
\end{extensionbox}

\begin{spiralbox}
\textbf{Coming next --- Lesson 2.5: Linear Regression.} We leave exact rules behind: given a
\emph{scatter} of real data points that don't lie on one line, we'll fit a single \textbf{line of best
fit} and measure how well it fits with the \textbf{correlation coefficient}, then use it to make
predictions.
\end{spiralbox}

\end{document}
